%created on August 18, 2005
%updated on december 19 2005


%\chapter {Problem Definition and Formulations}
\section{Problem Definition and Formulations}
\label{problem}

In this section, we provide a formal description of the LRSP 
considered in this research and present two formulations
of the problem.  

We consider an LRSP with capacity constraints on the facilities and 
on the vehicles, as well as time constraints on the vehicles.
In particular, given a set of candidate facility locations and a set of
customer locations, we seek a solution in which 
\begin{itemize}
\item[i.] each customer is visited exactly once, 
\item [ii.] each route starts and ends at the same facility, 
\item [iii.] the total demand of the customers assigned to a route is at 
most the vehicle capacity, 
\item [iv.] the total working time of a vehicle is no more than the  
time limit, and 
\item [v.] the total demand of the customers assigned to a facility 
does not exceed the capacity of the facility. 
\end{itemize}

We develop models for the delivery form of the LRSP and assume that there is no
service time at the customers.  The models, however, can be adapted easily for 
pickup problems and for nonzero service times. We assume that the demand 
information and distance matrix are known and not dynamic. 
We assume that each facility has enough vehicles to serve the allocated customers.
Simply, we assume that the number of vehicles at each facility is equal to the 
number of customers. And we consider one class of vehicles, all having the same 
characteristics and working under the same conditions. 

At the heart of the LRSP is a network flow structure, deriving from the routing 
of customer demands, that is constrained through the vehicle and facility 
capacities.  Hence, we consider both {\em arc-based} and {\em path-based} formulations 
for the LRSP, as each has been widely applied for related problems.
In Section \ref{edge-based-form}, we present a three-index commodity flow 
formulation of the LRSP. In Section \ref{set-part-form}, we present a 
set-partitioning-based formulation of the problem.  
%In Section \ref{sec:StrongForm},
%we introduce several classes of valid inequalities to strengthen the set-partitioning-based 
%formulation.

\begin{comment}
In this section, we present two formulations for a deterministic version 
of the LRSP. In Section \ref{edge-based-form}, we present two edge-based 
formulations of the problem. In Section \ref{set-part-form}, we present a set 
partitioning-based formulation of the problem.

Given a set of candidate facility locations and a set of customer locations, 
the objective of the LRS problem is to select a subset of the facilities, 
to construct a set of delivery routes and to assign routes to vehicles 
in such a way as to minimize total cost subject to the following 
constraints: (i) each customer must be visited exactly once, (ii) each 
route must start at a facility and return to the same facility, (iii) the 
total demand of the customers assigned to a route must be less than the vehicle 
capacity, (iv) the total working hours of a vehicle must be less than the working 
hour limit and (v) the total demand of the customers assigned to a facility 
must not exceed the capacity of the facility.

Throughout this section, we assume that the demand information and distance matrix 
are known and not dynamic. We consider the problem as a delivery 
problem, but the structure of the problem can easily be changed to a pickup problem. 
We assume that each facility has enough vehicles to serve the allocated customers.
Simply, we assume that the number of vehicles at each 
facility is equal to the number of customers. And we consider one class of vehicles, 
all having the same characteristics and working under the same conditions. 
\end{comment}

%\section{Edge-Based Formulations}
\subsection{Graph-Based Formulation of LRSP}
\label{edge-based-form}

For the related problems of LRP and VRP, the literature contains two main types of
graph-based models, vehicle flow and commodity flow. In both formulations, 
integer variables represent the number of times a vehicle traverses a given arc or 
edge in an underlying graph. In commodity flow formulations, an additional variable 
representing the number of units of a given commodity transported along the arc is
present. In general, vehicle flow models contain an exponential number of subtour 
elimination constraints, whereas commodity flow models, with the help of additional 
flow variables, use a polynomial number of constraints to eliminate subtours. 
For the LRP and the VRP, two-index vehicle flow formulations are preferred 
in solution algorithms based on branch-and-cut since these formulations include 
a smaller number of variables and their LP relaxations yield better bounds. 
However, it is not possible to formulate the LRSP using variables with only two 
indices because of the time constraint for the vehicles. In this context, we require a 
three-index commodity flow model.

To formulate the model, we introduce the following notation.

{\bf Sets}
\begin{eqnarray*}
I & = & \mbox{set of customer locations,} \\
J & = & \mbox{set of candidate facility locations,} \\
N & = & I \cup J = \mbox{set of all nodes,} \\
A & = & (J \times I) \cup (I \times I) = \mbox{set of arcs,}\\ 
H_j & = & \mbox{set of vehicles located at facility $j$, 
$\forall j \in J$, such that $|H_j|=|I|$,} \\ 
H & = & \bigcup_{j \in J}H_j = \mbox{set of all vehicles such that $|H|=|I||J|$.} 
\end{eqnarray*}


{\bf Parameters}
\begin{eqnarray*}
C^F_j & = & \mbox{daily equivalent fixed cost of opening facility $j$, $\forall j \in J$,} \\
C^O & = & \mbox{vehicle operating cost per unit travel time,} \\
C^V & = & \mbox{daily equivalent fixed cost of a vehicle (including the driver cost),} \\
D_i & = & \mbox{demand of customer $i$, $\forall i \in I$,} \\
L^F_j & = & \mbox{capacity of facility $j$, $\forall j \in J$,} \\
L^V & = & \mbox{capacity of a vehicle,} \\
L^T & = & \mbox{time limit for a vehicle and the driver, and} \\
T_{ij} & = & \mbox{travel time between locations $i$ and $j$, $\forall (i,j) \in A$.} 
\end{eqnarray*}

{\bf Decision Variables}
\begin{eqnarray*}
x_{ikh} & = & \left\{ \begin{array}{ll}
                 1  & \mbox{if vehicle $h$ travels on arc $(i,k)$, $\forall h \in H, (i,k) \in A$,} \\
                 0  & \mbox{otherwise,}                                                  
                \end{array} 
              \right. \\
y_{ikh} & = & \mbox{flow on arc $(i,k)$ carried by vehicle $h$, $\forall (i,k) \in A, h \in H$,} \\
t_{j} & = & \left\{ \begin{array}{ll}
                 1  & \mbox{if facility $j$ is selected, $\forall j \in J$,}\\
                 0  & \mbox{otherwise,}                                                  
                \end{array}  
              \right. \\
v_h & = &\left \{ \begin{array}{ll}
                 1  & \mbox{if vehicle $h$ is used, $\forall h \in H$, and}\\
                 0  & \mbox{otherwise.}                                                  
                \end{array}  \right. 
\end{eqnarray*}

%{\bf Formulation}
\begin{optprog}
{\bf (G-LRSP)} Minimize & \objective{\sum_{j \in J} C^F_j t_{j}  + C^V \sum_{h \in H} v_h  + C^O \sum_{h \in H} \sum_{(i,k) \in A}T_{ik} x_{ikh} } \label{obj-eb} \\
s.t. & \sum_{h\in H}\sum_{k \in N} x_{ikh} & = & 1 &  \forall  i \in I, \label{set-part}\\
           & \sum_{k\in N} x_{ikh}-\sum_{k \in N} x_{kih} & = & 0 & \forall i \in N, h \in H, \label{enterleave}\\
           & \sum_{h \in H_j}\sum_{k \in I} y_{jkh} - L^F_j t_{j} & \leq & 0 & 
	   \forall j \in J, \label{FacCap}\\
           & y_{ikh} - L^V x_{ikh} & \leq & 0 & \forall (i,k) \in A, h \in H, \label{VehCap}\\
           & \sum_{k \in N} y_{ikh} - \sum_{k \in N} y_{kih} + D_i\sum_{k \in N}x_{ikh} & = & 0 & \forall i \in I, h \in H, \label{flow-conv}\\
           & \sum_{(i,k) \in A} T_{ik} x_{ikh} - L^T v_h & \leq & 0 & \forall h \in H, \label{timeLim}\\
          % & \mbox{NumH} - \sum_{h \in H} \mbox{Dri}_h & = & 0 & \\
           %& \mbox{Tot} - \sum_{h \in H}\left[\sum_{i \in I}\sum_{k \in I}\left(t_{ik}+s_{k}\right)x_{ikh}\right] / 60 & = & 0 & \\
         %  & x_{jih} - T_{j} & \leq & 0 & \forall j \in M, i \in I, h \in H \label{Rel-X-T-1}\\
 %          & x_{ijh} - T_{j} & \leq & 0 & \forall j \in M ,i \in I, h \in H \label{Rel-X-T-2}\\
           & x_{jkh} & = & 0 & \forall j \in J,k \in N, h \in H_t, t \in J \backslash \{j\},\label{fix-0} \\ 
           & x_{ikh} & \in & \{0,1\} & \forall (i,k) \in A, h \in H, \label{x-var}\\
           & y_{ikh} & \geq &  0 & \forall (i,k) \in A, h \in H, \label{y-var} \\           
           & t_{j}  & \in & \{0,1\} & \forall j \in J, \mbox{ and} \label{t-var}\\
           & v_{h} & \in & \{0,1\} &  \forall h \in H. \label{dri-var}
        %  & \mbox{Tot} & \geq & 0 
\end{optprog}

The objective function (\ref{obj-eb}) states that the total cost, which  
includes the fixed cost of the selected facilities, the fixed cost of the vehicles
(including the driver cost), and the operating cost of the vehicles,
should be minimized. Constraints (\ref{set-part}) specify that exactly one vehicle 
must travel from customer node $i$ to some other node. Constraints (\ref{enterleave}) 
require that a vehicle should enter and leave a node an equal number of times. 
For customer nodes, constraints (\ref{set-part}) and (\ref{enterleave}) ensure 
that if a vehicle enters a 
customer node, it will leave this node and therefore that each customer 
node is served exactly once. For facility nodes, the number of visits by any 
vehicle can exceed one, since a vehicle can cover multiple routes. 
Constraints (\ref{FacCap}) ensure that the total outbound flows to the customer nodes 
from each facility do not exceed its capacity. 
Constraints (\ref{VehCap}) are vehicle capacity constraints and define the 
relationship between the binary variable $x_{ikh}$ and the flow variable 
$y_{ikh}$. Each constraint ensures that flow between any pair of nodes
cannot exceed the capacity of a vehicle. Constraints (\ref{flow-conv})
require conservation of flow at each customer node; these constraints
must be modified for pick-up problems.
Constraints (\ref{VehCap}) and (\ref{flow-conv}) 
together ensure that no route violates the vehicle capacity.
The combined sets of constraints (\ref{set-part}), (\ref{enterleave}),
and (\ref{flow-conv}) ensure that only valid traveling salesman
tours that include a  facility are formed, as in 
\citet {Nambiar}.  Constraints (\ref{timeLim}) limit the 
total time of a vehicle's schedule to the time limit. 
In this research, we assume that the set of vehicles in two different 
facilities are disjoint and a vehicle located at facility $j$ can only 
visit customer nodes and the facility $j$ during its trip.  
Constraints (\ref{fix-0}) restrict travel on the arcs originating from 
a facility to vehicles located at that facility.
Constraints (\ref{x-var}), (\ref{y-var}), (\ref{t-var}), and (\ref{dri-var}) 
are the integrality and non-negativity requirements on the variables.    

\begin{comment}
G-LRSP is a three-index commodity flow formulation.
Because of the facility and vehicle capacity restrictions, we cannot
know a priori the number of vehicles required at each facility, though the
number is clearly at most $|I|$.  Thus, this formulation can require
up to $O(|I|^4)$ binary variables and $O(|I|^4)$ constraints, which can be 
as many as 6 million for an instance with $50$ customers. 
Although it would be possible to develop a branch-and-bound algorithm for G-LRSP, 
we suspect that such an approach would not yield positive results, given the inherent
symmetry in the formulation and the known difficulty associated with solving related problems.
Instead, we propose a set-partitioning-based formulation 
similar to those that have been successful for other hard combinatorial problems.
\end{comment}

%%%%%%%%%%%%%%%%%%%%%%%%%%%%%%%%%%%%%%%%%%%%%%%%%%%%%%%%%%%%%%%%%%%%%%%
%%%%%%%%%%%%%%          COMMENT           %%%%%%%%%%%%%%%%%%%%%%%%%%%%
\begin{comment}
In model (EB-LRS), the objective function (\ref{obj-eb}) seeks to minimize the total cost, which  
includes the fixed cost of the selected facilities, the fixed cost of the 
drivers and vehicles and the operating cost of the vehicles. Constraints 
(\ref{set-part}) specify that exactly one vehicle must travel from demand node $i$
to some other node. Constraints (\ref{enterleave}) say that each vehicle should enter 
and leave a node an equal number of times. For demand nodes, constraints 
(\ref{set-part}) and (\ref{enterleave}) ensure that if a vehicle enters a 
demand node, eventually it leaves this node and therefore \zaupdate{that} each demand 
node is served exactly once. For facility nodes, the number of visits by any 
vehicle can be more than one, since a vehicle can cover multiple routes. 
Constraints (\ref{FacCap}) are capacity constraints that ensure that the 
capacity of a facility is not exceeded by flows to customer nodes. 
Constraints (\ref{VehCap}) are vehicle capacity constraints and define the 
relationship between the binary variable $x_{ikh}$ and flow variable 
$f_{ikh}$. Each constraint in (\ref{VehCap}) ensures that any flow between any 
node pair cannot be greater than the load capacity of a vehicle. Constraints (\ref{flow-conv})
require conservation of flow at each customer node based on the assumption that the 
problem is a delivery problem. This constraint should be changed if the type of
service is pickup instead of delivery. Constraints (\ref{VehCap}) and (\ref{flow-conv}) 
together ensure that the vehicle capacity is not violated in any route of the vehicle.
In combination with constraints (\ref{set-part}) and (\ref{enterleave}), 
constraints (\ref{flow-conv}) form the traveling salesman tours, as in 
\citet {Nambiar}.  Constraints (\ref{timeLim}) limit the 
total working time of a vehicle to the time limit. 
In model (EB-LRS), we assume that the set of vehicles in two different 
facilities are disjoint and a vehicle located at facility $j$ can only 
visit customer nodes and the facility $j$ during its trip.  
From this assumption, the model (EB-LRS) assumes that the $x$ and $f$ 
variables for the edges attached to a facility are set to zero (or have
large costs)
for vehicles that are not located at that facility. Constraints 
(\ref{set_to_0}) make sure that these variables are set to zero.
Constraints (\ref{x-var}), (\ref{t-var}), (\ref{dri-var}) and (\ref{y-var}) 
are the set of integrality and non-negativity requirements on the variables.    

%%%%%%%%%%%%%%%%%%%%%%%%%%%%%%%%%%%%%%%%%%%%%%%%%%%%%%%%%%%%%%%%%%%

The EB-LRS formulation is a three-index commodity flow model, and in 
the worst case,  with the assumption that $|J|\leq |I|$, it has $O(|I|^4)$ 
binary variables and $O(|I|^4)$ constraints. Therefore,
there are more than 6 million constraints and binary variables for an 
instance with $50$ customers. We could develop a branch-and-bound or
a branch-and-cut algorithm for the EB-LRS problem, however, we suspect
that such an approach will not yield good results, given the difficulty 
others have had in solving related problems with similar commodity formulations. 
For example, exact algorithms for the LRP generally are developed 
for two-index formulations of the problem, and none of the commodity 
flow models developed for the VRP is solved to optimality in the
literature. 

%Therefore, we instead are propose a set 
%partitioning-based approach, which has been successful for other hard 
%combinatorial problems. 
 
%Besides the fact that the edge-based formulation, EB-LRS has many binary variables 
%and many constraints, the EB-LRS is a three-index commodity flow model. 
%Exact algorithms for the LRP are generally developed for two-index formulations of the 
%problem.  Also, none of the commodity flow models developed for the VRP is solved to optimality in the
%literature. Furthermore, if there are side constraints for the vehicle routes such as 
%vehicle capacity, route length restriction and time windows, the decomposition algorithms 
%are common in the literature. Based on these observations, we develop an equivalent set 
%partitioning formulation for the LRS problem and we solve the set partitioning-based 
%formulation using a branch-and-price algorithm. 

%[moved to some other place]Although branch-and-price has been applied successfully to some of the 
%large scale optimization problems such as crew scheduling, machine
%scheduling, facility location and routing, since the mid-1990s, there are 
%many difficulties associated with the algorithm and the success of the 
%algorithm highly depends on the structure of the problem. \citet{Barnhart}
%provide an overview of the general methodology and many of the challenges.
%the formulation is not easily solved using conventional MIP methods
%%%%%%%%%%%%%%%%%%%%%%%%%%%%%%%%%%%%%%%%%%%%%%%%%%%%%%%%%%%%%%%%%%%%%%%%
\end{comment}


%%%%%%%%%%%%%%%%%%%%%%%%%%%%%%%%%%%%%%%%%%%%%%%%%%%%%%%%%%%%%%%%%%%%%%%
%%%%%%%%%%%%%%          COMMENT           %%%%%%%%%%%%%%%%%%%%%%%%%%%%

\begin{comment}
\subsection{A Two Index Vehicle Flow Formulation}
By removing the flow variable $f$ from (EB-LRS), rephrasing facility capacity constraint 
(\ref{FacCap}) with $x$ variables and replacing vehicle capacity (\ref{VehCap}) 
and flow conservation constraints (\ref{flow-conv}) with rounded capacity inequalities
for each vehicle, we can obtain a three index vehicle flow formulation of LRSP. 
\citet{Belenguer06} present a two index vehicle flow formulation for a location
routing problem (LRP) with capacitated facilities and capacitated vehicles which is a
special case of LRSP. To obtain a two index vehicle flow formulation of the
LRSP, we need to add a scheduling constraint to the LRP to guarantee that the
time limit is not exceeded while assigning routes to vehicles. 
In distance constraint VRP (DCVRP), \citet{Laporte5} use optimal cost of multiple TSP 
to drive a lower bound for the minimum number of vehicles based on the distance 
constraint for vehicle routes. The time constraint in the LRSP is different from the
distance constraint in the DCVRP in the sense that the time constraint in LRSP is for a 
vehicle not for a single route and it is possible to assign multiple routes to a vehicle.

To formulate the LRSP with two index variables, we duplicate each facility and
add $|J|$ additional nodes representing the copies of the facilities to
the problem. Let $J^A$ be the set of artificial facilities and $a(j) \in J^A$ be the
index of the facility corresponds to facility $j \in J$. Let $V$ be the set of 
nodes, i.e., $V = I \cup J \cup J^A$. We consider the symmetric version   
of the problem (i.e., $T^T_{ik}=T^T_{ki}$ for all $i,k \in V$) and thus a undirected
network. Let $E$ be the set of edges. There are not edges between any facility and a
copy of another facility. 
   
{\bf Decision Variables}
\begin{eqnarray*}
x_{e} & = & \left\{ \begin{array}{ll}
                 1  & \mbox{if edge $e$ is in the solution, $\forall e \in E$}\\
                 0  & \mbox{otherwise}                                                  
                \end{array} 
              \right. \\
t_{j} & = & \left\{ \begin{array}{ll}
                1  & \mbox{if facility $j$ is selected to be open, $\forall j \in J$}\\
                 0  & \mbox{otherwise}                                                  
                \end{array}  
              \right. 
\end{eqnarray*}
To simplify the notation, let $\delta(i)$ be the set of edges connected to node $i \in V$, 
$\delta(S)$ denote the set of edges with one end in $S$ and the other in $V \backslash S$
for all $S \subseteq V$, $\gamma(S)$ be the set of edges with both end in $S$ and 
$D^C(S)=\sum_{i \in S}D^C_i$ for all $S \subseteq I$.  

Using the same parameters in section \ref{3-ind-CFform}: 
\begin{optprog}
{\bf (EB-LRS2)} Min & \objective{\sum_{j\in J} C^F_j t_{j}  + C^V/2 \sum_{j \in J}\sum_{e \in \delta(j)}x_{e}  + C^O \sum_{e \in E}T^T_{e}x_{e}} \label{obj-eb2} \\
s.t. & \sum_{e \in \delta(i)} x_{e} & = & 2 &  \forall  i \in I \label{set-part2}\\
           & x_{e}-2y_{j} & \leq & 0 & \forall j \in J, e \in \delta(j) \cup \delta(a(j)) \label{con3}\\
           & \sum_{e \in \delta(S\cup \{j\})} x_{e}  & \geq & 2 & 
	   \forall j \in J, \forall S \subseteq I | D^C(S) > L^F_j  \label{FacCap2}\\
           & \sum_{e \in \delta(S)} x_{e}  & \geq & 2 \lceil D^C(S) / L^V \rceil & \forall S \subseteq I \label{VehCap2}\\
	   & \sum_{e \in \delta(S)} x_{e}  & \geq & 2 \lceil B^C(S) / L^T \rceil & \forall S \subseteq I \cup J^A \label{timeLim2}\\
     
%      & \sum_{e\in \gamma(S)\cup \{i\}}x_{e} + x_{ji}-  \sum_{e\in \gamma(S)\cup \{i\}}& \leq & 0 & \forall h \in H \label{timeLim}\\


	%   & \sum_{i\in V} x_{jih} & = & 0 & \forall j \in J, h \in H \backslash H_j \label{set_to_0}\\

           & x_{e} & \in & \{0,1,2\} & \forall j \in J \cup J^A, e \in \delta(j)  \label{x-var2}\\
	   & x_{e} & \in & \{0,1\} & \forall S \subseteq I | |S|=2, e \in \gamma(S)  \label{x-var3}\\
	   & t_{j}  & \in & \{0,1\} & \forall j \in J \label{t-var2}
\end{optprog}
\end{comment}


%\section{A Set Partitioning Based Formulation}
\subsection{Set-Partitioning-Based Formulation of LRSP}
\label{set-part-form}
In this section, we introduce a set-partitioning-based formulation of 
the LRSP. First, we define the column set constructed to formulate
the LRSP. 
 
\subsubsection{Pairing Definition}
To define the column set, the set-partitioning-based formulation 
uses the idea of a {\em pairing}, which we adapt from the pairing concept in the 
crew scheduling literature (e.g., \citet{Desrosiers}, \citet{Vance}, and \citet{Anbil}).
In the crew scheduling literature, a pairing is a sequence of duty periods 
and represents  a possible schedule for a single crew. In the LRSP context, 
a pairing is a sequence of routes that represents a possible (feasible) schedule for a 
single vehicle.
%In the crew scheduling literature, a pairing is a sequence of duty periods and represents  
%a possible schedule for a single crew. In the LRS context, a pairing is a 
%set of routes that can be served sequentially by one vehicle within the 
%vehicle's working hour limit. 

A pairing is feasible if 
\begin{itemize}
\item [i.] each route included in the pairing starts and ends at the 
same facility, 
\item [ii.]for each route in the pairing, the total demand of the 
customers assigned to the route is at most the vehicle capacity, 
\item [iii.] the total travel time of the pairing is at most the vehicle time limit,
\item [iv.] each customer included in the pairing must be visited only once.
\end{itemize}
By adding constraint (iv), we prevent the generation of pairings that
visit some customer(s) multiple times since such pairings are set 
to zero in an integer feasible solution of an LRSP.

Figure \ref {fig:exppair} shows two possible pairing constructions.  
Pairing 1 includes routes 1, 2 and 3; pairing 2 includes routes 4 and 5. 
Each pairing specifies the routes covered by a vehicle and has an associated  
cost, which is defined as the sum of the costs of the component routes plus a
fixed cost for the vehicle. Since the cost of a multiple-stop route depends on the sequence 
of nodes visited, the pairing information implicitly includes sequencing information
for each route. The cost of a pairing, however, is not affected by the ordering of
the routes in the absence of time windows and thus is arbitrary.  We will take this
fact into consideration to improve the solution algorithms for the subproblem and 
will not generate symmetric pairings corresponding to different orderings of the same routes.

\begin{figure}[h]
  \centering
  \includegraphics[scale=0.5]{figs/pairexp2.eps}
  %\fbox{ \includegraphics[scale=0.5]{figure/pair2.eps}}
  \caption{Example of pairings} \label{fig:exppair}
\end{figure}

\subsubsection{Formulation}
\label{part1:sp_form}
Given a set of candidate facility locations and the set of feasible pairings that can 
be constructed for each candidate facility, the objective of the LRSP is to 
select a subset of the facilities and an associated subset of pairings in such a way 
as to minimize the total cost subject to the following constraints: 
\begin{itemize}
\item [i.] each customer must be covered by exactly one pairing, 
\item [ii.] the total demand of the customers assigned to a facility must not 
exceed the capacity of the facility, and  
\item [iii.]  vehicles cannot be associated with a facility unless it is open. 
\end{itemize}
To present the formulation, we define additional parameters and 
decision variables.  Associated with each candidate facility 
$j \in J$, we let $P_j$ be the set of all feasible pairings for facility
$j$.   

{\bf Parameters}
\begin{eqnarray*}
a_{ip} & = & \left \{  \begin{array}{ll}
                 1 & \mbox{if customer node $i$ is in pairing $p$ of facility $j$, $\forall i \in I, p \in P_j, j \in J$,} \\
                 0 & \mbox{otherwise, and}
                         \end{array}
              \right. \\
C_p & = & \mbox{cost (vehicle fixed cost plus operating cost) of pairing $p$ %associated with 
of facility $j$, $\forall p \in P_j, j \in J$.} 
\end{eqnarray*}
{\bf Decision Variables}
\begin{eqnarray*}
t_j &  = & \left \{ \begin{array}{ll}
                 1  & \mbox{if facility $j$ is selected, $\forall j \in J$,}\\
                 0  & \mbox{otherwise,}                                                  
                \end{array}
       \right. \\
z_p & = & \left \{ \begin{array}{ll}
                 1  & \mbox{if pairing $p$ is selected for facility $j$, $\forall p \in P_{j}, j \in J$, and}\\
                 0  & \mbox{otherwise.}                                                  
                \end{array}
       \right.
\end{eqnarray*}
\begin{optprog}
{\bf(SPP-LRSP)} Minimize & \objective{\sum_{j \in J} C^F_j t_{j} + \sum_{j \in J} \sum_{p\in P_j} C_p z_p} \label{spp-obj}\\
subject to & \sum_{j \in J} \sum_{p \in P_{j}} a_{ip} z_{p} & = & 1 & \forall i \in I, \label{spp_con1} \\
           & \sum_{p \in P_j} \sum_{i \in I} a_{ip} D_i z_{p} - L^F_j t_j & \leq & 0 & \forall j \in J, \label{spp_con2} \\
           & t_{j}  & \in & \{0,1\} & \forall  j \in J, \mbox{ and}\label{binary1} \\
           & z_{p} &  \in & \{0,1\} & \forall p \in P_j, j \in J. \label{binary2} 
\end{optprog} 
 

The objective function (\ref{spp-obj}) seeks to minimize the total cost, which 
includes the fixed cost of the selected facilities and the cost of the selected pairings.
We note that the cost $C_{p}$ of a pairing includes both the vehicle fixed cost and 
variable costs. Constraints (\ref{spp_con1}) 
guarantee that each customer location is served by exactly one pairing. 
%states that each customer must be in exactly one of the selected pairings.
%Constraints (\ref{spp_con2}) guarantee that the total demand of the customers assigned to
%a facility does not exceed the facility capacity.  
Constraints (\ref{spp_con2}) ensure that the total demand of the selected pairings does 
not exceed the facility capacity.  Constraints (\ref{binary1}) and (\ref{binary2}) are 
standard binary restrictions on the variables. 


%Constraints (\ref{vehicle-bound}) set an upper bound for the total number of pairings 
%(i.e., number of vehicles) that can be selected for each facility. Based on our assumption in
%section (\ref{edge-based-form}), $|H_j|=|I|$, since the cost of a pairing includes fixed cost
%of vehicles and the size of set $|H_j|$ is large enough, we can drop this constraint,
%it is redundant. 
%Constraints (\ref{binary1}) and (\ref{binary2}) are standard binary restrictions on the 
%variables. 
The structure of the SPP-LRS is very similar to the CFLP with single sourcing (CFLPSS). In the
CFLPSS, only pairings with one customer are considered and assigned to capacitated
facilities. However, in the LRSP, pairings with multiple customers must be taken into account. 


\section{Strengthening the Set-Partitioning-Based Formulation}
\label{StrongForm}
In this section, we discuss how we can strengthen the set-partitioning-based formulation
SPP-LRSP by adding valid inequalities. First, we introduce a lower bound on the total 
number of facilities required in any integer feasible solution.  Let $R$ denote a lower bound
on the minimum number of required facilities. 
Then we can add to SPP-LRSP a constraint to force the number of selected
facilities to be at least $R$:
\begin{eqnarray}
\sum_{j \in J} t_{j} \geq R. \label{lb-open-fac}
\end{eqnarray}
For any instance, we can calculate a valid value for $R$ from the demands of the customers and
the capacities of the facilities. There are two possible cases:\\
\noindent {\bf Case 1.} If all of the facilities have the same capacity value, 
$L^F_j = k$ $\forall j \in J$, then 
\begin{equation}
R =\left \lceil \frac{\sum_{i \in I} D_i}{k} \right \rceil . \nonumber
\end{equation}
\noindent {\bf Case 2.} If at least one facility capacity differs, then a lower bound 
on the number of required facilities can be calculated using a greedy algorithm: 
compute the total demand and assign it to the facilities in non-increasing order of their 
capacities.  Let $n = |J|$ and let $(j_1,j_2,..,j_n)$ be an ordering of the facilities such that 
$L^F_{j_1}\geq L^F_{j_2} \geq ...\geq L^F_{j_n}$. Then,
\begin{eqnarray*}
 R = \mbox{argmin}_{\{l = 1..n\}} \left (\sum_{s=1}^{l}L^F_{j_s}\geq \sum_{i \in I} D_i \right ).  
\end{eqnarray*}

Constraint (\ref{lb-open-fac}) is not necessary in the formulation of the LRSP. However, as a simple
lower bound on the number of required facilities, it may eliminate some solutions in which 
a fractional number of facilities is used, thereby improving the LP relaxation bound.
 The following proposition can be stated:

\begin {prop} 
Chv\'atal-Gomory rank (CG rank) of constraint (\ref{lb-open-fac}) is 1, if
$\left (\frac{\sum_{i \in I }D_i}{k}\right)$ is fractional in case 1 and
$\left (\frac{\sum_{i \in I }D_i}{L^F_{j_1}}\right)$ is fractional or strictly less 
than $R$ in case 2. 
\end {prop}

\begin{proof}

Recall that the CG rank of any original inequality and any 
inequality dominated by a nonnegative linear combination of the original inequalities 
is 0.  An inequality which can be generated as a CG inequality based on rank 0 inequalities 
but is not rank 0 itself has CG rank 1. 

\noindent {\bf Case 1:} $L^F_j = k$ $\forall j \in J$.

To prove the proposition, we compute the summation of constraints (\ref{spp_con2}) over all 
facilities $j \in J$:
\begin{eqnarray}
 \sum_{j \in J} \sum_{p \in P_j} \sum_{i \in I} a_{ip} D_i z_{p} & \leq & k \sum_{j \in J} t_{j} \label{sum-spp_con2}
\end{eqnarray} 
Then, if we substitute constraint (\ref{spp_con1}) in (\ref{sum-spp_con2}) and divide by $k$, we get:
\begin{eqnarray}
 \sum_{j \in J} t_{j} &\geq &\frac {\sum_{i \in I}D_i}{k} \label{sum2-spp_con2}
\end{eqnarray} 

The $t_j$ variables are binary, so $\sum_{j \in J}t_{j}$ must be integer. 
Therefore, if we let 
$R =\left \lceil \frac{\sum_{i \in I }D_i}{k} \right \rceil$, then constraint
(\ref{lb-open-fac}) must be valid. If  the right hand side of constraint (\ref{sum2-spp_con2})
%$\frac {\sum_{i \in I}D_i}{k}$
is fractional, the constraint (\ref{lb-open-fac}) has CG rank 1. Otherwise, its rank is 0,
since it can be defined as the non-negative linear combination of the original inequalities.
 

\noindent{\bf Case 2:} The facilities have different capacity values. 

To generalize this statement for the problems in which the facility
capacities might be different in value, we reorder the facilities as
$t_{[1]},t_{[2]},..,t_{[J]}$ such that $L^F_{[1]} \geq L^F_{[2]} \geq 
..\geq L^F_{[J]}$. Then, inequality (\ref{sum-spp_con2}) becomes 
\begin{eqnarray}
\sum_{j \in J} \sum_{p \in P_j} \sum_{i \in I} a_{ip} D_i z_{p} \leq   \sum_{j \in J} L^F_{[j]}t_{[j]} \label{nonhomo-sum}
\end{eqnarray}  
Again, if we substitute in (\ref{nonhomo-sum}) using constraint (\ref{spp_con1}) and 
then divide each side by $L^F_{[1]}$, 
we get:
\begin{eqnarray}
\label{nonhomo-sum2}
 t_{[1]}+\frac{L^F_{[2]}}{L^F_{[1]}}t_{[2]}+..+ \frac{L^F_{[J]}}{L^F_{[1]}}t_{[J]} \geq \frac {\sum_{i \in I}D_i}{L^F_{[1]}} 
\end{eqnarray} 
If we let 
\begin{eqnarray}
 R_2 =\left \lceil \frac{\sum_{i \in I }D_i}{L^F_{[1]}} \right \rceil \nonumber 
\end{eqnarray}   
then, by rounding up the coefficients of inequality (\ref{nonhomo-sum2}), we get:
\begin{eqnarray}
\label{lb-open-fac2}
 \sum_{j \in J}t_j \geq R_2 
\end{eqnarray}  

The value of $R_2$ is calculated based on the assumption that all facilities have the 
same capacity, which is the largest capacity value. However, in the calculation procedure 
of $R$ for problems that have different capacity values (case 2), it is considered that the 
second facility might not have a capacity as large as the first facility, the third facility 
might not have a capacity as large as the second facility and so on.  Therefore, 
$R \geq R_2$ and the constraint (\ref{lb-open-fac}) dominates constraint 
(\ref{lb-open-fac2}). If $\left(\frac{\sum_{i \in I }D_i}{L^F_{[1]}}\right)$ is fractional or 
$R_2<R$, then the CG rank of constraint (\ref{lb-open-fac}) is 1. \qed
\end{proof}

Due to the similarity in structure between the SPP-LRSP and the CFLP, we might consider 
adding the set of constraints 
\begin{eqnarray}
z_{p} \leq & t_{j}  & \forall j \in J, \forall p \in P_{j} \label{relation-Z-T}
\end{eqnarray}
In the CFLP, constraints similar to inequality (\ref{relation-Z-T}) are facets of the 
CFL polytope if the demands of the customers in the constraints are less than the capacity 
of the facilities (\citet{Leung}). In addition, \citet{Daskin} remarks that constraints 
in the CFLP similar to (\ref{relation-Z-T}) considerably improve the LP bound. 

Constraints (\ref{relation-Z-T})  ensure that a pairing of a facility can be selected 
only if that facility is selected. Although constraints (\ref{relation-Z-T}) are implied 
by constraints (\ref{spp_con2}), constraints (\ref{relation-Z-T}) may improve the 
LP relaxation of the model as in the case of CFLP. However, since there is potentially 
an exponential number of such constraints, we would like to find a smaller set of 
constraints to add. 
\begin {prop} 
  To strengthen the LP relaxation bound, we can replace constraints (\ref{relation-Z-T}) with 
  \begin{eqnarray}
    \label{relation}
    \sum_{p \in P_{j}}a_{ip}z_{p}\ \leq \ t_{j} \mbox{\ \ $\forall j \in J$ and $i \in I$} 
  \end{eqnarray}
and constraints (\ref{relation}) have a CG rank of 1.
\end{prop}
 
\begin{proof} 
Constraints (\ref{relation}) are valid for the integer program. In any feasible solution,
each demand node $i \in I$ is assigned to exactly one facility, so 
$\sum_{p \in P_j} a_{ip} z_{p}$ is $1$ for one facility and $0$ for all others. 
For any $j \in J$ such that $\sum_{p \in P_j} a_{ip} z_{p}$ is $1$, $t_j$ will be $1$, 
which corresponds to pairings being selected only for open facilities. 

%In any feasible integer solution, a demand node 
%$i \in N$ can be assigned to only one facility, so 
%$\sum_{p \in P_j} a_{ip} Z_{p}$ can only be 1 or 0 for any facility $j \in M$. Also 
%constraints (\ref{spp_con2}) guarantee that if any demand node is assigned to facility j,
%then $T_j$ must be greater than zero. Since $T_j$ is binary variable, if $T_j>0$, $T_j$ must be 1.
%Therefore, the inequality (\ref{relation}) is valid for any feasible integer solution. 

Also, when the coefficients of the variables in inequality (\ref{relation}) and constraint 
(\ref{relation-Z-T}) are compared, it is easy to see that inequality (\ref{relation}) 
dominates constraint (\ref{relation-Z-T}). For example, if we take $z_{p'} \leq t_{j'}$
from constraint set (\ref{relation-Z-T}) for $j=j'$ and $p=p'$, we can compare it with a
constraint in set (\ref{relation}) for which $j=j'$ and $i=i'$ such that $a_{i'p'}=1$. 
That is $z_{p'}+\sum_{p \in P_{j'}\setminus\{p'\}}a_{i'p}z_{p}\ \leq \ t_{j'}$. Because 
of the $z_{p}$ terms other than  $z_{p'}$ on the left side of the constraint,   
$\sum_{p \in P_{j'}}a_{i'p}z_{p}\ \leq \ t_{j'}$ dominates $z_{p'} \leq t_{j'}$. 

Now we show the rank of the inequality $\sum_{p \in P_j} a_{ip} z_{p} \leq T_j$ 
for an arbitrary $i$ and $j$. For notational convenience, suppose $i=1$ and $j=1$.  

Let $D(k)=\sum_{i \in I}a_{ik}D_i$ be the total demand of a pairing $ k\in P_1$ . 
From constraint sets (\ref{spp_con1}) ($\leq$ part of the constraints) and (\ref{spp_con2}), 
we can write:
\begin{eqnarray}
 \sum_{\{p \in P_1 | a_{1p}=1\}} z_{p} + \sum_{\{j \in J \setminus \{1\}\}}\sum_{\{p \in P_{j}|a_{1p}=1\}}z_{p}  &\leq & 1 \label{prop2-a} \\
 \sum_{\{p \in P_1 | a_{1p}=1\}} D(p) z_{p} + \sum_{\{p \in P_1 | a_{1p}=0\}} D(p) z_{p} - L^F_1 t_1 & \leq & 0 \label{prop2-b} 
\end{eqnarray} 

Let $D_{min} = min_{\{p \in P_1 | a_{1p}=1\}} \{ D(p) \}$. Multiply inequality (\ref{prop2-a}) 
with $(L^F_1-D_{min})$ and add the resulting inequality to the inequality (\ref{prop2-b}). 
Then after dividing the obtained inequality by $L^F_1$ round it to floor: 

\begin{eqnarray}
 \sum_{\{p \in P_1 | a_{1p}=1\}}  \left \lfloor \frac{L^F_1-D_{min}+D(p)}{L^F_1} \right \rfloor z_{p} &+&
    \sum_{\{j \in J / \{1\}\}}\sum_{\{p \in P_{j}|a_{1p}=1\}}\left \lfloor \frac{L^F_1-D_{min}}{L^F_1}\right \rfloor z_{p}\nonumber \\
     &+&\sum_{\{p \in P_1 | a_{1p}=0\}}\left \lfloor \frac{D(p)}{L^F_1}\right \rfloor z_{p} - t_1\nonumber \\
&\leq&  \left \lfloor \frac {L^F_1-D_{min}}{L^F_1}\right \rfloor  \label{prop2-c}
\end{eqnarray} 


Also, we simply assume that $D(p) \leq L^F_1 \ \ \forall p \in P_1$ since otherwise we can set
$z_{p}=0$ and remove the variable from the constraint. This is true, because if the
total demand of the pairing is greater than the capacity of the facility, this pairing
cannot be chosen in a feasible solution for this facility. In inequality (\ref{prop2-c}), 
the coefficients in the first summation are in the range $[1,2)$, the coefficients in the second 
summation are in $[0,1)$, and the coefficients in the third summation are in the range $[0,1)$ if 
$D(p)<L^F_1$ for all $p \in P_1$. The right hand side is in the range $[0,1)$. 
Thus when we round down the coefficients, we get $\sum_{p \in P_1} a_{1p} z_{p} \leq t_1$. 
Since we have obtained inequality (\ref{relation}) by using the original constraints and the 
CG inequality procedure, the CG rank of the inequality is 1. Also,  this inequality is not 
dominated by any linear combination of the original constraints, thus its CG rank cannot be 0. 
If $D(p)=L^F_1$ for some $p \in P_1$ in the third summation, when we round down, 
the coefficient for this variable in (\ref{prop2-c}) is 1. Thus, we can add the variables $z_p$ 
for all $p \in P_1$ and satisfying $D(p)=L^F_1$ and  $a_{1p}=0$ to inequality (\ref{relation}) 
to get a stronger constraint. But if we do not add these variables, the constraint (\ref{relation}) 
will still be valid. 
\qed
\end{proof}

